\documentclass[letterpaper,11pt]{article}

% T1: acentos como glifo único (í, no ı+´) y versalitas extraíbles como texto.
% Sin esto, un ATS o un Ctrl+F por "clínicos" no hace match.
\usepackage[T1]{fontenc}
\usepackage[empty]{fullpage}
\usepackage[spanish]{babel}
\usepackage{titlesec}
\usepackage[hidelinks]{hyperref}
\usepackage{enumitem}
\usepackage{tabularx}
\usepackage{charter}


% ==================== MÁRGENES ====================

\addtolength{\oddsidemargin}{-0.5in}
\addtolength{\evensidemargin}{-0.5in}
\addtolength{\textwidth}{1in}
\addtolength{\topmargin}{-0.5in}
\addtolength{\textheight}{1.3in}


% ==================== FORMATO DE SECCIONES ====================

% Mayúsculas reales en vez de \scshape: Charter no trae versalitas propias,
% las simula escalando la misma fuente y pdftotext lee el cambio de tamaño
% como espacio ("E XPERIENCIA"). \MakeUppercase extrae limpio.
\titleformat{\section}
  {\vspace{-4pt}\raggedright\large}
  {}
  {0em}
  {\MakeUppercase}
  [\titlerule \vspace{-5pt}]


% ==================== COMANDOS PROPIOS ====================

\newcommand{\resumeItem}[1]{%
  \item\small{#1 \vspace{-2pt}}%
}

\newcommand{\resumeSubheading}[4]{%
  \vspace{-2pt}\item
  \begin{tabular*}{0.97\textwidth}[t]{l@{\extracolsep{\fill}}r}
    \textbf{#1} & #2 \\
    \textit{\small#3} & \textit{\small #4} \\
  \end{tabular*}\vspace{-7pt}%
}

\newcommand{\resumeProjectHeading}[2]{%
  \item
  \begin{tabular*}{0.97\textwidth}{l@{\extracolsep{\fill}}r}
    \small#1 & #2 \\
  \end{tabular*}\vspace{-7pt}%
}

\newcommand{\resumeSubHeadingListStart}{%
  \begin{itemize}[leftmargin=0.15in, label={}]%
}
\newcommand{\resumeSubHeadingListEnd}{%
  \end{itemize}\vspace{-3pt}%
}

\newcommand{\resumeItemListStart}{\begin{itemize}}
\newcommand{\resumeItemListEnd}{\end{itemize}\vspace{-3pt}}

\raggedright
\raggedbottom
\setlength{\tabcolsep}{0in}


% ==================== DOCUMENTO ====================

\begin{document}


% -------------------- ENCABEZADO --------------------

\begin{center}
  \textbf{\Huge Felipe Farias Escobar} \\ \vspace{1pt}
  \small Ingeniero de Software Semi Senior \\ \vspace{1pt}
  Santiago, Chile \\
  \href{mailto:felipe.farias.e1@gmail.com}{felipe.farias.e1@gmail.com} $|$
  \href{https://www.linkedin.com/in/FelipeFa6/}{Linkedin.com/in/FelipeFa6} $|$
  \href{https://github.com/FelipeFa6}{Github.com/FelipeFa6}
\end{center}\vspace{-6pt}


% -------------------- RESUMEN PROFESIONAL --------------------

\section{Resumen Profesional}

\small{
  Ingeniero de Software con 4 años de experiencia en desarrollo full stack,
  en producción con \textbf{TypeScript} sobre \textbf{React} y
  \textbf{NestJS}, además de \textbf{Go} y \textbf{Python}, en los sectores
  de salud, energía y retail.
}


% -------------------- EXPERIENCIA PROFESIONAL --------------------

\section{Experiencia Profesional}

\resumeSubHeadingListStart

  \resumeSubheading
    {Ingeniero de Software Semi Senior}{Junio 2025 -- Julio 2026}
    {INDRA}{Híbrido}
  \resumeItemListStart

    \resumeItem{Desarrollo de frontend en \textbf{React} con
    \textbf{TypeScript} para plataformas del \textbf{Coordinador Eléctrico
    Nacional}: vistas y componentes nuevos, integración con APIs, refactor,
    migración de código existente y corrección de bugs.}

    \resumeItem{Desarrollo sobre microservicios en \textbf{NestJS} con
    TypeScript: endpoints y funcionalidades nuevas, corrección de bugs y
    ejecución del despliegue. Mantenimiento de microservicios en \textbf{Go},
    desarrollo con \textbf{Django} sobre \textbf{AWS} y despliegues sobre
    \textbf{GCP} (Compute Engine).}

    \resumeItem{Implementación de la solución de Calidad del Dato junto al
    arquitecto de datos: reglas de validación y listeners \textbf{event saga}
    (Java/Quarkus, Python). Calidad de datos de \textbf{75\% a 98\%},
    validada externamente. Participación en arquitectura, \textbf{code
    review} y optimización de queries.}

  \resumeItemListEnd

  \resumeSubheading
    {Ingeniero de Software}{Enero 2024 -- Junio 2025}
    {Instituto Nacional del Cáncer}{Presencial}
  \resumeItemListStart

    \resumeItem{Propuesta, diseño y construcción desde cero del
    \textbf{Integrador}, orquestador concurrente en \textbf{Go} que sincroniza
    \textbf{4 bases de datos Oracle} de sistemas clínicos aislados mediante
    endpoints HTTP propios, con patrón \textbf{event saga}, compensación ante
    fallos parciales y reintentos persistentes. \textbf{Sigue en producción}.}

    \resumeItem{Interoperabilidad \textbf{FHIR} vía \textbf{Mirth Connect}.
    \textbf{CI/CD} con Git Hooks y migración a \textbf{Docker} (despliegue de
    \textbf{1 hora a 5 minutos}); legacy \textbf{PHP 5.7/CodeIgniter 3}
    migrado a \textbf{PHP 8/Laravel}.}

  \resumeItemListEnd

  \resumeSubheading
    {Desarrollador de Software Junior}{Julio 2022 -- Diciembre 2023}
    {Coolebra}{Remoto}
  \resumeItemListStart

    \resumeItem{Actualización completa de la landing page en \textbf{React}:
    sección de productos con descuento y variación de precios que resolvía
    búsquedas contra \textbf{OpenSearch}, panel de descuentos por tarjetas y
    visualización de evolución de precios (\textbf{Chart.js}) para verificar
    si un descuento estaba inflado. \textbf{Mercado Pago} para
    suscripciones.}

    \resumeItem{Construcción desde cero de un \textbf{ETL} en \textbf{Python}
    sin dependencias externas, que auditaba \textbf{200 mil productos
    diarios} desde \textbf{PostgreSQL} hacia OpenSearch, en \textbf{2 horas}
    sobre \textbf{5 máquinas} (\textbf{Terraform}), y escalamiento del
    scraping de \textbf{más de 30 retailers} sobre \textbf{3 instancias EC2}.}

  \resumeItemListEnd

  \resumeSubheading
    {Tutor de Programación Orientada a Objetos III y Bases de Datos}
    {Noviembre 2021}
    {INACAP}{Presencial}
  \resumeItemListStart

    \resumeItem{Docencia en herencia y modelado orientado a objetos en
    \textbf{Python} y bases de datos relacionales.}

  \resumeItemListEnd

\resumeSubHeadingListEnd


% -------------------- PROYECTO ACTUAL --------------------

\section{Proyecto Actual}

\resumeSubHeadingListStart

  \resumeProjectHeading
    {\textbf{Petfi.io} $|$ Plataforma social de rescate y adopción}
    {Marzo 2026 -- Presente}
  \resumeItemListStart

    \resumeItem{Plataforma con más de \textbf{11 mil usuarios activos}.
    Aplicación web en \textbf{Flutter Web}, incluido el panel de
    administración y su \textbf{sistema de roles}, sobre \textbf{Firebase}
    (Firestore, Auth, Storage, Messaging, Hosting) con reglas de seguridad e
    índices compuestos. Scripts en \textbf{Node.js} y \textbf{Spec Driven
    Development}.}

  \resumeItemListEnd

\resumeSubHeadingListEnd


% -------------------- EDUCACIÓN Y CERTIFICACIONES --------------------

\section{Educación y Certificaciones}

\resumeSubHeadingListStart

  \resumeProjectHeading
    {\textbf{INACAP} $|$ Ingeniería en Informática $|$ Titulado 2023}
    {2019 -- 2022}

  \resumeProjectHeading
    {\textbf{University of Helsinki} $|$ Full Stack Open, curso principal
    (React, Redux, Node.js, Express, MongoDB)}
    {2024}

  \resumeProjectHeading
    {\textbf{Oracle} $|$ OCI 2025 Certified Foundations Associate}
    {2026}

\resumeSubHeadingListEnd


% -------------------- TECNOLOGÍAS --------------------

\section{Tecnologías}

\small{
  \textbf{Lenguajes}: TypeScript, JavaScript, Go, Python, SQL, PHP, Dart \\
  \textbf{Frontend}: React, Chart.js, Flutter Web \\
  \textbf{Backend}: NestJS, Go/Gin, Django, FastAPI, Laravel \\
  \textbf{Datos}: PostgreSQL, Oracle/PL-SQL, OpenSearch, Firestore, MongoDB \\
  \textbf{Infraestructura y Prácticas}: AWS (EC2), GCP (Firebase, Compute
  Engine), Docker, Terraform, CI/CD, Linux, microservicios, Event Saga, ETL,
  testing, code review, Scrum
}


\end{document}
